\documentclass[UTF8,fontset=windows,12pt]{ctexart}
%\setcounter{tocdepth}{1}
%\setcounter{secnumdepth}{4}
%\usepackage{ctex} %若要输入中文请取消注释或者把article改成ctexart
% This first part of the file is called the PREAMBLE. It includes
% customizations and command definitions. The preamble is everything
% between \documentclass and \begin{document}.
%\usepackage[OT1]{fontenc}
%\usepackage{fontspec}
%\setmainfont{Times New Roman}
%\setCJKmainfont{Microsoft YaHei}
\usepackage[left=1.91cm,right=1.91cm,top=2.54cm,bottom=2.54cm]{geometry}
\usepackage{amsfonts}
\usepackage{amsmath}
\usepackage{amssymb}
\usepackage[upint,lcgreekalpha]{stix}
\usepackage[type1]{garamondlibre}
%\usepackage{esint}
\usepackage{pgfplots}
\pgfplotsset{compat=1.18}
\usepackage{tikz}
\usetikzlibrary{arrows}
\usepackage{extarrows}
\usepackage[only,llbracket,rrbracket]{stmaryrd}
\usepackage{titlesec}
\usepackage{bm}
\usepackage{graphicx}
\usepackage{appendix}
\usepackage{tikz-cd}
\usepackage{tikz-3dplot}
\usetikzlibrary{arrows.meta, decorations.pathreplacing, calc, patterns}
\usepackage{float}
\usepackage{mathtools}
\usepackage{amsthm}
\usepackage{verbatim}              % better theorem environments
\usepackage{setspace}
\usepackage{enumitem}
\setenumerate[1]{itemsep=0pt,partopsep=0pt,parsep=\parskip,topsep=5pt}
\setitemize[1]{itemsep=0pt,partopsep=0pt,parsep=\parskip,topsep=5pt}
\setdescription{itemsep=0pt,partopsep=0pt,parsep=\parskip,topsep=5pt}
\usepackage{xcolor}
\usepackage[colorlinks,linkcolor=black,anchorcolor=black,citecolor=black]{hyperref}
%\usepackage{apacite}
\usepackage[numbers]{natbib}
%\usepackage{showkeys}
\usepackage{cases}
\usepackage{fancyhdr}
\usepackage[scaled=0.92]{helvet}	% set Helvetica as the sans-serif font
%\usepackage{upgreek}
%\renewcommand{\rmdefault}{ptm}		% set Times as the default text font
%\usepackage{newtxtext}
\bibliographystyle{plain}
% If AMS-LaTeX is used, it can be loaded before or after mtpro2		
% The following loads metro and defines some common MTPro options [2, 4]
%\usepackage[lite,subscriptcorrection,nofontinfo]{mtpro2}
%\pdfmapfile{+mtpro2.map}

%\usepackage{txfonts}

% various theorems, numbered by section

\makeatletter
\renewenvironment{proof}[1][\proofname]{\par
  \pushQED{\qed}%
  \normalfont \topsep6\p@\@plus6\p@\relax
  \trivlist
  \item[\hskip\labelsep
        \bfseries % 关键修改：斜体改为加粗
        #1\@addpunct{.}]\ignorespaces
}{%
  \popQED\endtrivlist\@endpefalse
}
\makeatother
\newtheoremstyle{kaiti-theorem}   % 样式名称
  {3pt}                           % 上方间距
  {3pt}                           % 下方间距
  {\kaishu\upshape}               % 正文字体：中文楷体 + 英文直立
  {}                              % 缩进量
  {\bfseries\upshape}      % 头部字体（定理名）：中文楷体粗体 + 英文直立粗体
  {.}                             % 定理名后标点
  {0.5em}                         % 定理名后间距
  {}                              % 头部说明（留空）

% 应用自定义样式
\theoremstyle{kaiti-theorem}
\newtheorem{thm}{定理}[section]
\newtheorem{nota}[thm]{记号}
\newtheorem{question}{问题}
\newtheorem*{questionn}{思考}
\newtheorem{exe}{习题}[section]
\newtheorem{assump}[thm]{假设}
\newtheorem*{thmm}{定理}
\newtheorem{defn}{定义}[section]
\newtheorem{lem}[thm]{引理}
\newtheorem*{lemm}{引理}
\newtheorem{prop}[thm]{命题}
\newtheorem*{propp}{命题}
\newtheorem*{claim}{断言}
\newtheorem{cor}[thm]{推论}
\newtheorem{conj}[thm]{猜想}
\newtheorem{rmk}{注记}[section]
\newtheorem{ex}{例}[section]
\newtheorem{pf}{证明}
\newtheorem{problem}{作业题}
\numberwithin{equation}{section}





\DeclareMathOperator{\id}{id}
\DeclareMathOperator*{\esssup}{ess\,sup}
\renewcommand{\Re}{\textbf{Re}\,}  
\renewcommand{\Im}{\textbf{Im}\,}  
\newcommand{\R}{\mathbb{R}}  
\newcommand{\FH}{\mathbb{H}}  
\newcommand{\C}{\mathbb{C}}  
\newcommand{\Z}{\mathbb{Z}}
\newcommand{\X}{\mathbb{X}}
\newcommand{\N}{\mathbb{N}}
\newcommand{\Q}{\mathbb{Q}}
\newcommand{\T}{\mathbb{T}}
\newcommand{\TT}{\mathcal{T}}
\newcommand{\TP}{\overline{\partial}{}}
\newcommand{\TJ}{\langle\TP\rangle}
\newcommand{\TL}{\overline{\Delta}{}}
\newcommand{\curl}{\text{curl }}
\newcommand{\dive}{\text{div }}
\newcommand{\bd}[1]{\mathbf{#1}}  % for bolding symbols
\newcommand{\bds}[1]{\boldsymbol{#1}}  % for bolding symbols
\newcommand{\RR}{\mathcal{R}}      % for Real numbers
\newcommand{\sh}{\mathcal{S}}  
\newcommand{\sph}{\mathbb{S}}  
\newcommand{\Om}{\Omega}
\newcommand{\OmT}{{\Omega_T}}
\newcommand{\ut}{{U_T}}
\newcommand{\vp}{\varphi}
\newcommand{\Th}{\Theta}
\newcommand{\Lam}{\Lambda}
\newcommand{\Omb}{{\overline{\Omega}}}
\newcommand{\OmbT}{{\overline{\Omega_T}}}
\newcommand{\ub}{{\overline{U}}}
\newcommand{\ubt}{{\overline{U_T}}}
\newcommand{\q}{\quad}
\newcommand{\p}{\partial}
\newcommand{\pb}{\boldsymbol{\partial}}
\newcommand{\lap}{\Delta}
\newcommand{\dd}{\mathfrak{D}}
\newcommand{\DD}{\mathcal{D}}
\newcommand{\den}{{\bar\rho}_{0}}
\newcommand{\Dt}{{\p}_{t}+v^{k}{\p}_{k}}
\newcommand{\col}[1]{\left[\begin{matrix} #1 \end{matrix} \right]}
\newcommand{\noi}{\noindent}
\newcommand{\nab}{\nabla}
\newcommand{\cnab}{\overline{\nab}}
\newcommand{\cp}{\overline{\partial}{}}
\newcommand{\dx}{\,\mathrm{d}x}
\newcommand{\dxx}{\,\mathrm{d}\bds{x}}
\newcommand{\xxi}{\boldsymbol{\xi}}
\newcommand{\eeta}{\boldsymbol{\eta}}
\newcommand{\dxxi}{\,\mathrm{d}\boldsymbol{\xi}}
\newcommand{\deeta}{\,\mathrm{d}\boldsymbol{\eta}}
\newcommand{\dxi}{\,\mathrm{d}\xi}
\newcommand{\deta}{\,\mathrm{d}\eta}
\newcommand{\dy}{\,\mathrm{d}y}
\newcommand{\dyy}{\,\mathrm{d}\bds{y}}
\newcommand{\dz}{\,\mathrm{d}z}
\newcommand{\dzz}{\,\mathrm{d}\bds{z}}
\newcommand{\dph}{\,\mathrm{d}\varphi}
\newcommand{\dt}{\,\mathrm{d}t}
\newcommand{\dr}{\,\mathrm{d}r}
\newcommand{\dtt}{\,\mathrm{d}\tau}
\newcommand{\dS}{\,\mathrm{d}S}
\newcommand{\dss}{\,\mathrm{d}\sigma}
\newcommand{\dmu}{\,\mathrm{d}\mu}
\newcommand{\dnu}{\,\mathrm{d}\nu}
\newcommand{\dV}{\,\mathrm{d}V}
\newcommand{\ds}{\,\mathrm{d}s}
\newcommand{\drr}{\,\mathrm{d}\rho}
\newcommand{\dist}{\text{dist }}
\newcommand{\vol}{\text{vol}\,}
\newcommand{\volo}{\text{vol}(\Omega)}
\newcommand{\spt}{\text{Spt}\,}
\newcommand{\Tr}{\text{Tr}\,}
\newcommand{\lee}{\langle}
\newcommand{\ree}{\rangle}
\newcommand{\xee}{\langle\xi\rangle}
\newcommand{\xxee}{\langle\boldsymbol{\xi}\rangle}
\newcommand{\xeta}{\langle\eta\rangle}
\newcommand{\xeeta}{\langle\boldsymbol{\eta}\rangle}
\newcommand{\kk}{\kappa}
\newcommand{\lam}{\lambda}
\newcommand{\io}{\int_{\Omega}}
\newcommand{\iu}{\int_{U}}
\newcommand{\is}{\int_{\Sigma}}
\newcommand{\ipo}{\int_{\partial\Omega}}
\newcommand{\ipu}{\int_{\partial U}}
\newcommand{\ig}{\int_{\Gamma}}
\newcommand{\ir}{\int_{\R}}
\newcommand{\ird}{\int_{\R^d}}
\newcommand{\irn}{\int_{\R^n}}
\newcommand{\fpi}{\frac{1}{(2\pi)^{\frac{d}{2}}}}
\newcommand{\ddt}{\frac{\mathrm{d}}{\mathrm{d}t}}
\newcommand{\ddx}{\frac{\mathrm{d}}{\mathrm{d}x}}
\newcommand{\eps}{\varepsilon}
\providecommand{\jump}[1]{\left\llbracket #1 \right\rrbracket }
\providecommand{\len}[1]{\lee #1 \ree }
\providecommand{\ino}[1]{\left\| #1 \right\| }
\providecommand{\bno}[1]{\left| #1 \right| }
\newcommand{\red}{\textcolor{red}}
\newcommand{\green}{\textcolor{green}}
\newcommand{\blue}{\textcolor{blue}}
\newcommand{\nnr}{{[n+1]}}
\newcommand{\nnn}{{[n]}}
\newcommand{\nnl}{{[n-1]}}
\newcommand{\nnll}{{[n-2]}}	
\newcommand{\mmr}{{[m+1]}}
\newcommand{\mmm}{{[m]}}
\newcommand{\mml}{{[m-1]}}
\newcommand{\mmll}{{[m-2]}}

%---------------Special variables--------------
\newcommand{\CC}{\mathfrak{C}}  
\newcommand{\NN}{\mathbf{N}}
\newcommand{\LH}{\mathcal{L}}
\newcommand{\HH}{\mathcal{H}}  
\newcommand{\EE}{\mathcal{E}}  
\newcommand{\FT}{\mathcal{F}}  
\newcommand{\fT}{\mathbf{T}}  
\newcommand{\FA}{\mathbf{A}}  
\newcommand{\MH}{\mathcal{M}}
\newcommand{\NH}{\mathcal{N}}
\newcommand{\PH}{\mathcal{P}}
\newcommand{\VH}{\mathcal{V}}
\newcommand{\XX}{\mathfrak{X}}
\newcommand{\xx}{{\bds{x}}}
\newcommand{\bx}{\mathbf{x}}
\newcommand{\by}{\mathbf{y}}
\newcommand{\bp}{\mathbf{p}}
\newcommand{\bq}{\mathbf{q}}
\newcommand{\pp}{{\bds{p}}}
\newcommand{\qq}{{\bds{q}}}
\newcommand{\uu}{\mathbf{u}}
\newcommand{\ww}{\mathbf{w}}
\newcommand{\vv}{\mathbf{v}}
\newcommand{\yy}{{\bds{y}}}
\newcommand{\zz}{{\bds{z}}}
\newcommand{\zo}{\mathbf{0}}
\newcommand{\ee}{\mathbf{e}}
\newcommand{\fa}{\mathbf{a}}
\newcommand{\fb}{\mathbf{b}}
\newcommand{\fc}{\mathbf{c}}
\newcommand{\ff}{\bds{f}}
\newcommand{\fr}{\mathbf{r}}
\newcommand{\fh}{\mathbf{h}}
\newcommand{\fm}{\mathbf{m}}
\newcommand{\fg}{\mathbf{g}}
\newcommand{\FF}{\mathbf{F}}
\newcommand{\GG}{\mathbf{G}}
\newcommand{\BB}{\mathbf{B}}



\begin{document}
\title{\LARGE{\bf  2026年秋季学期~常微分方程~作业二}}
\author{解的延拓、比较原理、解对初值和参数的连续依赖性}
\date{截止时间：2026.9.28下课前（电子版截至当天23:59:59）}

\maketitle

\textbf{禁止抄袭或使用AI直接生成答案PDF，被抓到者此次作业得零分。}

%考虑如下形式的一阶常微分方程的初值问题
%\begin{equation}\label{2-0: IVP}\tag{2.0.1}
% \xx'(t)=\ff(t, \xx(t)), \qquad \xx(t_0)=\xx_0, 
%\end{equation}
%其中 \(I\subset\R\) 是含有 \(t_0\) 的区间, 未知函数\(\xx:I\to\R^d\), 给定函数 \(\ff:D\to\R^d\) 的定义域\(D\) 是 \(\R\times\R^d\) 中的开集, 并且 \((t_0, \xx_0)\in D\). 

\begin{problem}[习题2.3.3]
设初值问题 $x'=(x^2-2x-3)e^{(x+t)^2},~ x(t_0)=x_0$ 的解的最大存在区间为 $(a,b)$. 证明: $a=-\infty$ 与 $b=+\infty$ 中至少有一个成立.

提示：先找出这个方程的所有平衡态（即不随$t$变化而变化的解），然后讨论初值$x_0$与这个平衡态之间的关系。
\end{problem}

\begin{problem}[习题2.3.4] 
证明初值问题 $x'=t^3-x^3,~ x(t_0)=x_0$ 的解在 $[t_0,+\infty)$ 上存在。

提示：为了使用延拓准则，可以考虑去估计 $x(t)^2$ 这个量的演化，最终可以算得 $\ddt(x(t)^2)\leqslant \frac32 |t|^4$。
\end{problem}

%\begin{problem}[问题2.3.2]
%设 $f$ 在 $\R$ 上是连续且单调递增的，$C$ 是常数。考虑初值问题 $x'(t)=-f(x(t)) + C$, $x(t_0)=x_0$.
%\begin{itemize}
%\item [(1)] 证明：对任意给定的常数 $C$, 该初值问题在 $[t_0,+\infty)$ 上存在解。
%\item [(2)] 该方程在$[t_0,+\infty)$上的解是否唯一，证明你的结论。
%\end{itemize}
%
%提示：设时间正向的极大存在区间是 $[t_0,\beta)$, $f$ 在 $\R$ 上的下确界、上确界为 $m, M$ (可能为无穷). 分情况讨论 $C>M$, $C<m$, $m\leqslant C\leqslant M$. 对最后一种情况, \underline{若}存在 $a$ 使得 $f(a)=C$, 则讨论 $|x(t)-a|^2$ 的单调性.
%\end{problem}

\begin{problem}[问题2.3.3]
设 $f(t,x)$ 在 $0\leqslant t\leqslant a,\ -\infty<x<+\infty$ 上连续，且存在常数 $q\in (0,1)$ 使得对任意的 $t\in (0,a]$ 都有 $ |f(t,y)-f(t,z)| \leqslant \dfrac{q}{t}|y-z|.$ 证明：初值问题 $x'=f(t,x),\ x(0)=x_0$ 在 $[0,a]$上存在唯一解。

提示：在证明存在性时，Peano存在性定理只给出某个小区间 $[0,\delta]$ 上的存在性，思考如何将其延拓到整个 $[0,a]$ 上；在证明唯一性时，思考为什么我们要求$q\in(0,1)$.
\end{problem}


\begin{problem}[问题2.4.1，选做]
考虑初值问题 $x'(t)=f(t,x)$, $x(t_0)=x_0$, 其中 $f(t,x)$ 在 $\R^2$ 连续并关于 \(x\) 局部 Lipschitz。设 $\alpha,\beta\in C^1([t_0,+\infty))$ 满足
\[
 \alpha(t)\leqslant\beta(t),\qquad \alpha'(t)\leqslant f(t,\alpha(t)), \qquad \beta'(t)\geqslant f(t,\beta(t)).
\]
\begin{enumerate}
 \item [(1)] 证明 $x_0\in[\alpha(t_0),\beta(t_0)]$ 时对应的解满足对任意 $t\geqslant t_0$ 都有 $\alpha(t)\leqslant x(t)\leqslant\beta(t)$
 \item [(2)] 证明该解正向整体存在。
 \item [(3)] 若$x_0\in(\alpha(t_0),\beta(t_0))$, 条件中的所有不等式改为严格不等式, 则(1)的结论也变为严格不等式.
\end{enumerate}

提示：(1) 对 $\alpha_\eps(t):=\alpha(t)-\eps e^{K(t-t_0)}$ 和 $x(t)$ 本身用比较原理，然后取极限 $\eps\to 0_+$ 得到不等式左边，右边可构造类似的辅助函数。
\end{problem}

\begin{problem}[问题2.5.1]
考虑 $[t_0,\infty)$ 上的含参初值问题 $\xx'(t)=f(t,\xx(t), \lam)$, $\xx(t_0)=\xx_0$. 其中 $\ff: [t_0,+\infty)\times\R^d\times\R^m\to \R^d$ 连续，并关于 $\xx$ 局部 Lipschitz.  假设存在常数
\(\alpha>0,~L\geqslant0,~C\geqslant0\), 满足
\begin{align*}
(\xx-\yy)\cdot(\ff(t,\xx,\lam)-\ff(t,\yy,\lam))&\leqslant -\alpha|\xx-\yy|^2,\\
 |\ff(t,\xx,\lam)-\ff(t,\xx,\mu)|\leqslant L|\lam-\mu|,&\qquad |\ff(t,\zo,\lam)| \leqslant C(1+|\lam|).
\end{align*}
\begin{enumerate}
 \item [(1)] 证明：对每个固定的参数 $\lam\in\R^m$, 该初值问题的解都在时间正向整体存在。
 \item [(2)] 记 $\varphi(t;t_0,\xx_0,\lam)$ 是以 $\lam$ 为参数, $t_0$ 时刻初值为 $\xx_0$ 的解，证明
 \[
  |\varphi(t;t_0,\xx_0,\lam)- \varphi(t;t_0,\yy_0,\mu)|\leqslant e^{-\alpha(t-t_0)}|\xx_0-\yy_0|+ \frac{L}{\alpha}\bigl(1- e^{-\alpha(t-t_0)}\bigr) |\lam-\mu|.
 \] 并说明解对初值和参数的依赖为何关于 $t\in[t_0,\infty)$ 是一致的。
\end{enumerate}

提示：(1) 考察 $|\xx(t)|^2$ 的演化，(2)  如果令 $\ww:=\xx-\yy$, 则有 $\ddt|\ww|\leqslant -\alpha|\ww|+L|\lam-\mu|$. 严格来说，如果要避免$|\ww|$在0处不可微，则应当先考虑其 $C^1$ 逼近 $\ww_\eps(t):=\sqrt{\eps^2+|\ww(t)|^2}$，得到它的估计之后再令 $\eps\to 0$.  请注意这是一种对绝对值函数常用的$C^1$逼近。
\end{problem}

\begin{problem}[习题2.6.1]
对含参微分方程初值问题 $x'(t)=\lam x(t)-x(t)^3$, $x(0)=x_0$ 写出关于 $x_0$ 和 $\lam$ 的变分方程, 并证明 $\dfrac{\p x(t)}{\p x_0}>0.$
\end{problem}




\begin{problem}[习题2.6.4]
考虑微分方程 $x'(t)=\cos^2x(t) + (\lam+q(t))\sin^2 x(t)$, $x(0;\lam)=x_0$, 其中 $q\in C(\R)$.
\begin{itemize}
\item [(1)] 证明：对任意 $\lam\in\R$, 该初值问题在 $t\in[0,1]$ 上存在唯一解 $x=\vp(t;\lam)$.
\item [(2)] 对任意 $t\in[0,1]$, 这个解$x=\vp(t;\lam)$是否关于 $\lam\in\R$ 是连续可微的, 证明你的结论.
\item [(3)] 令 $\omega(\lam)=\vp(1;\lam)$, 证明: $\omega(\lam)$ 关于 $\lam$ 严格单调递增.
\end{itemize}

提示：(3) 使用关于参数 $\lam$ 的变分方程，请注意我们需要说明 $\frac{\p\vp}{\p\lam}(1;\lam)$ 是\underline{严格}正的。
\end{problem}


\end{document}
